\documentclass[11pt]{article}
\usepackage[margin=1in]{geometry}
\usepackage{amsmath,amssymb}
\title{Exact Transfer Recurrence and Phase Root Count for $n=4$}
\author{Wendel Le\'on Alvarado}
\date{November 12, 2026}
\begin{document}
\maketitle

\begin{abstract}
We give an exact finite-state frontier transfer for the $n=4$ Ehrhart
polynomials and prove its cubic recurrence for every layer parameter.
The squared factor $q_6^2$ in the characteristic discriminant creates
two eigenvalue contacts inside the audited interval. Tracking the
analytic eigenvalue branch gives a phase change of $-3\pi$, while its
spectral coefficient cannot be continued as a single continuous branch
through both contacts. Splitting at these contacts and applying
80-digit interval bounds gives $k$ strict sign changes in each of three
windows for every $k\geq3$. Degree saturation, together with exact
Sturm verification of the first actual case $k=2$, proves that every
polynomial associated with the order-polytope family ($k\geq2$) has
exactly $3k$ distinct, real, negative roots. The $k=1$ polynomial used
in the recurrence is an auxiliary seed, not an Ehrhart polynomial of
this family.
\end{abstract}

\section{Exact frontier transfer}
The order-polytope family considered here is indexed by $k\geq2$.
To state the recurrence from its first transfer step, define the
auxiliary sequence
\[
 \widehat H_{0,4}=1,\qquad
 \widehat H_{1,4}=t^3+11t^2+11t+1,\qquad
 \widehat H_{k,4}=H_{k,4}\quad(k\geq2).
\]
In particular, $\widehat H_{1,4}$ is a recurrence seed and is not
$h^*(\mathcal O(C_{1,4}))$.
Let $h=(h_0,h_1,h_2,h_3)$ describe an order ideal at a cut with $q$
complete layers: column $i$ contains its first $h_i$ elements. The
constraints are
\[
  \sum_i h_i=4q,\qquad 0\leq h_i\leq q+1,\qquad
  h_i\leq h_{i-1}+1
\]
with cyclic indices. Put $d_i=h_i-q$. Then $\sum_i d_i=0$ and $d_i\leq
1$, so $d_i\geq-3$. For every $q\geq3$, the admissible deviation vectors
therefore lie in the same fixed finite set
\[
  d_i\in\{-3,-2,-1,0,1\},\qquad
  \sum_i d_i=0,\qquad d_i\leq d_{i-1}+1.
\]
A maximal element in column $c$ is its top element and exists precisely
when $h_c>0$ and $h_{c+1}\leq h_c$. The canonical state records the
occupancy masks of rows $q-1$ and $q$, together with the last element's
row offset $d_c$ and column $c$. Exact enumeration gives 20 states.
The canonical state sets at $q=2$ and $q=3$ also agree, so the same
one-layer transition applies from the base window onward.

Adding one layer uses only these two boundary rows. The descent weight
depends on the previous element's relative row and column, not on $q$.
Thus the transition matrix $M(t)$ is independent of the number of
complete layers. Prefix and suffix vectors $v(t)$ and $u(t)$, obtained
from the three-layer base window, give
\[
  H_{k,4}(t)=u(t)M(t)^{k-3}v(t),\qquad k\geq3.
\]
All arithmetic here is in $\mathbb Z[t]$. The exact 20-state computation
verifies
\[
  \bigl(M^3-P_4M^2-Q_4M-R_4I\bigr)v=0,
\]
where
\[
\begin{aligned}
P_4(t)&=t^3+23t^2+23t+1,\\
Q_4(t)&=7t^5-28t^4-22t^3-28t^2+7t,\\
R_4(t)&=t(t-1)^6(t+1).
\end{aligned}
\]
Multiplying the vector identity by $uM^j$ proves the scalar recurrence
for all $k\geq6$; exact checks at $k=3,4,5$ supply the initial cases.
Consequently,
\[
  \widehat H_{k,4}
  =P_4\widehat H_{k-1,4}
   +Q_4\widehat H_{k-2,4}
   +R_4\widehat H_{k-3,4}
  \qquad(k\geq3).
\]
The transfer also reproduces the exact certificate polynomials
$H_{3,4}$, $H_{4,4}$, and $H_{5,4}$. The finite dynamic-programming
data independently verifies the same identity for $3\leq k\leq12$.
The vector annihilator is sufficient for this scalar recurrence; we do
not claim that the full $20\times20$ characteristic polynomial equals
the displayed cubic.

The proposed second-order relation with constant boundary term
\[
H_{k,4}-P_4H_{k-1,4}-Q_4H_{k-2,4}=R_4
\]
is false. The exact residuals for $k=4$ and $k=5$ are, respectively,
\[
\begin{aligned}
k=4:\;&t^2(t-1)^6(t+1)(t^2+11t+11),\\
k=5:\;&t^2(t-1)^6(t+1)\cdot
(t^5+38t^4+263t^3+484t^2+263t+38).
\end{aligned}
\]

\section{Phase audit and root-count consequence}
The audited interval is
\[
 I_4=(\beta_-,\beta_+),\qquad
 \beta_\pm=\frac{z_0\pm\sqrt{z_0^2-4}}2,\qquad
 z_0=-38-24\sqrt3.
\]
It is reciprocal-symmetric, contains $t=-1$, and covers both sides of
$-1$. The phase calculation uses the eigenvalue branch
\(\widetilde\lambda_2\) with positive imaginary part. The discriminant
factors exactly as
\[
\Delta_X(t)=-4tq_4(t)q_6(t)^2,
\]
where
\[
 q_4=t^4+76t^3-282t^2+76t+1,\qquad
 q_6=t^6+7t^5+27t^4+58t^3+27t^2+7t+1.
\]
There are exactly two simple roots of $q_6$ in $I_4$,
\[
 \tau_1=-3.4372313754626764\ldots,\qquad
 \tau_2=-0.2909318258697067\ldots,
\]
with signs $q_6:+,-,+$ on the three successive windows. The roots are
isolated by rational intervals in the machine-readable certificate.
Only $t=-1$ is a root of $R_4=t(t-1)^6(t+1)$ in $I_4$.

Put $p=-Q_4-P_4^2/3$, $q=-2P_4^3/27-P_4Q_4/3-R_4$, and let $s(t)$
be the simple real root of $s^3+ps+q=0$. Writing
$D=3s^2+p>0$, exact reduction modulo this cubic gives
\[
 \left(\frac{3s^2}{4}+p\right)D^2=tq_4q_6^2.
\]
Thus the analytic eigenvalue branch initialized in the upper half-plane
has
\[
 \widetilde\lambda_2(t)
 =-\frac{s}{2}+\frac{P_4}{3}
   +i\,q_6(t)\frac{\sqrt{tq_4(t)}}{D}.
\]
The signed imaginary part is allowed to change sign at both $q_6$ roots.
At $\tau_1$ the double eigenvalue is positive and the branch crosses
the positive real axis from above to below. At $\tau_2$ the double
eigenvalue is negative and the branch crosses the negative real axis
from below to above. The latter crossing requires subtracting $2\pi$
from the principal argument. Consequently the continuous analytic
eigenvalue lift has
\[
 \Delta\arg\widetilde\lambda_2=-3\pi.
\]
The old $-\pi$ value is the winding of the reflected,
positive-imaginary-part branch. No factor $W(t)$ or normalization is
responsible for the missing turn.

There is an important obstruction to assigning a global continuous
phase to the coefficient of this analytic eigenvalue branch. At the
two contacts the interval-enclosed one-sided upper-branch coefficient
limits are,
respectively,
\[
 \begin{aligned}
 \widetilde c_2^+(\tau_1)
  &=0.34665821575\ldots+0.12174458208\ldots i,\\
 \widetilde c_2^+(\tau_2)
  &=0.34665821575\ldots-0.12174458208\ldots i.
 \end{aligned}
\]
Their arguments have absolute value
$\theta_1,\theta_2\in(0,\pi/4)$, certified with 80-digit interval
arithmetic. Following the analytic eigenvalue through a contact
requires conjugating its coefficient, so the one-sided coefficient
limits are not equal. To check the confluent representation, set
$D_6=P_4^2+3Q_4$, $\nu_d=-(P_4Q_4+9R_4)/2$, and
$\nu_s=P_4D_6-2\nu_d$. Exact reduction gives
\[
 D_6^2H_{3,4}-D_6(\nu_s+\nu_d)H_{2,4}
       +\nu_s\nu_d \widehat H_{1,4}\equiv0\pmod{q_6}.
\]
Since the repeated and simple eigenvalues are nonzero and distinct at
these contacts, this identity removes the generalized $k\lambda^k$
term. The confluent pair formula is therefore valid there, but there
is no single global continuous $\arg\widetilde c_2$ to combine with the
$-3\pi$ eigenvalue lift.

For root counting, use the upper-half-plane eigenvalue and coefficient
separately on each of the three open windows. Their limiting phases give
\[
\begin{array}{c|c|c|c}
 \text{endpoint}&\arg\widetilde\lambda_2^+
 &\arg\widetilde c_2^+&\Phi_k^+\\ \hline
 \beta_-&\pi&\pi/2&(k+\tfrac12)\pi\\
 \tau_1&0&\theta_1&\theta_1\\
 \tau_2&\pi&-\theta_2&k\pi-\theta_2\\
 \beta_+&0&-\pi/2&-\pi/2
\end{array}
\qquad
\Phi_k^+=\arg\widetilde c_2^+
                 +k\arg\widetilde\lambda_2^+.
\]
The exact coefficient formulas are continuous on each window, and the
two contact values are their certified one-sided limits. The
2,000-cell, 80-digit interval audit proves
\[
 \mathcal E_2(t)<0.02,\qquad
 \frac{|\lambda_1(t)|}{|\lambda_2(t)|}<0.24
 \quad(t\in I_4).
\]
Hence for every $k\geq2$,
$\mathcal E_k\leq\mathcal E_2<0.02<1/\sqrt2$. At phase levels that
are integer multiples of $\pi$, or at a contact phase $\theta_i$ or
$k\pi-\theta_i$, the dominant pair therefore determines the strict
sign of $H_{k,4}$.

For $k\geq3$, the following ordered sample levels lie in the indicated
windows and have alternating signs:
\[
\begin{array}{ll}
(\beta_-,\tau_1):&
 k\pi,(k-1)\pi,\ldots,\pi,\theta_1,\\
(\tau_1,\tau_2):&
 \theta_1,\pi,2\pi,\ldots,(k-1)\pi,k\pi-\theta_2,\\
(\tau_2,\beta_+):&
 k\pi-\theta_2,(k-1)\pi,\ldots,\pi,0.
\end{array}
\]
The phase is continuous on each window and its endpoint limits bracket
the listed levels, so the intermediate value theorem supplies ordered
sample points even without a monotonicity claim. Their strict
alternating signs give at least $k$ sign changes in each window.
The recurrence has degree $3k$ with leading coefficient $1$ (the
$P_4H_{k-1,4}$ term supplies the leading term), so these $3k$ distinct
negative roots exhaust the degree. The actual first case $H_{2,4}$ has
six negative roots by exact Sturm. Thus every Ehrhart polynomial
$H_{k,4}$ in the stated family has exactly $3k$ distinct negative roots
for $k\geq2$. Separately, exact Sturm verifies that the auxiliary
recurrence seed $\widehat H_{1,4}$ has three distinct negative roots;
this is not a statement about a $k=1$ order polytope.
Independent exact Sturm checks for $k=3,\ldots,12$ place all $3k$ roots
in $[-79,-1/79]\subset I_4$.

The 80-digit derivative enclosures at the two $q_6$ contacts are
\[
\begin{array}{c|r|r|r|l}
 &\lambda_{\rm simple}&\lambda_{\rm double}
 &\lambda_{\rm simple}'&\widetilde\lambda_2'\\ \hline
 \tau_1&
 13.04455168&70.01231057&-10.09722209&
 -44.78587130-23.18722179i\\
 \tau_2&
 -0.32122021&-1.72404306&0.37472125&
 4.74816937+6.74590077i
\end{array}
\]
At $\beta_-$ and $\beta_+$, the repeated eigenvalue pair has unbounded
derivatives because $q_4$ has a simple zero there; the remaining
eigenvalue derivatives are enclosed near $19823.30482837$ and
$269.20437641$, respectively. At $t=-1$, the eigenvalues are
$0,\pm4\sqrt3\,i$, with derivatives $-4/3$ and
$-28/3\mp6\sqrt3\,i$. All rational isolators and full interval
enclosures are recorded in
\texttt{certificado\_teorema\_n4.json}.

\section{Spectrum and roots in the right exterior window}
Write $\alpha=-\beta_+$, so
$\alpha=0.012569659525542648674850600415\ldots$. The characteristic
discriminant factors exactly as
\[
 \Delta_X(t)=-4t\,q_4(t)\,q_6(t)^2,
\]
where
\[
q_4=t^4+76t^3-282t^2+76t+1,\qquad
q_6=t^6+7t^5+27t^4+58t^3+27t^2+7t+1.
\]
There is exactly one $q_4$ zero in $(-1,0)$, namely $-\alpha$.
There is also one $q_6$ zero in $(-1,0)$, isolated in
$(-1/3,-1/4)$; it lies strictly inside $I_4$, to the left of
$-\alpha$. Hence there is no discriminant zero in $(-\alpha,0)$.
Since $t<0$ and $q_4(t)>0$ there, $\Delta_X(t)>0$: all three
eigenvalues are real and simple throughout this open window.
Indeed, with $z=t+t^{-1}$, the two factors reduce to
$q_4/t^2=z^2+76z-284$ and
$q_6/t^3=z^3+7z^2+24z+44$; the exact Sturm isolations select the
negative roots described above.

An equal-modulus pair of distinct real eigenvalues must be opposite.
That requires
\[
 P_4Q_4+R_4=
 8t(t+1)(t^2-5t+1)(t^4+20t^3+22t^2+20t+1)=0.
\]
The negative roots of this expression are outside $(-\alpha,0)$.
Thus no eigenvalue-modulus crossing occurs in the right window. At the
rational sample $t=-1/1000$, 80-digit interval enclosures isolate the
eigenvalues in increasing order near
$-0.02831489035$, $0.03664098649$, and $0.96869690286$.
The last is the unique dominant positive eigenvalue, and its coefficient
is enclosed near $1.02552505798$. Continuity and the absence of modulus
ties keep this positive real branch dominant throughout $(-\alpha,0)$.

Using $\widehat H_{0,4},\widehat H_{1,4},H_{2,4}$, write
\[
 c_j=\frac{N(\lambda_j)}{\partial_X\chi_4(\lambda_j)},\qquad
 N(X)=X^2\widehat H_{0,4}
       +X(\widehat H_{1,4}-P_4\widehat H_{0,4})
       +(H_{2,4}-P_4\widehat H_{1,4}-Q_4\widehat H_{0,4}).
\]
The exact resultant is
\[
 \operatorname{Res}_X(\chi_4,N)
   =16t^2(t-1)^2q_6(t)^2.
\]
It has no zero in $(-\alpha,0)$; consequently none of the three
simple-mode coefficients vanishes there, and the dominant coefficient
$c_1$ stays positive. There is therefore no second complex-dominance
window and no alternating pair with equal modulus in $(-\alpha,0)$.

Exact Sturm counts for recurrence-generated $H_{k,4}$ give:
\[
\begin{array}{c|rrrrrrrrrr}
k&3&4&5&6&7&8&9&10&11&12\\ \hline
\#\text{ roots in }I_4&9&12&15&18&21&24&27&30&33&36\\
\#\text{ roots in }(-\alpha,0)&0&0&0&0&0&0&0&0&0&0\\
\#\text{ roots in }(-\infty,-1/\alpha)&0&0&0&0&0&0&0&0&0&0
\end{array}
\]
Indeed, all $3k$ roots are counted in the rational interval
$[-79,-1/79]\subset I_4$, so degree saturation leaves no roots in
either exterior window. The proposed $(k,k,k)$ split is therefore
already contradicted by exact data: for $k=3,\ldots,12$ the distribution
is $(3k,0,0)$, not $(k,k,k)$.

This exterior-window analysis is consistent with, but is not needed for,
the all-$k$ root count: the three-window argument above already places
all $3k$ roots inside $I_4$. The finite Sturm table additionally shows
that for $k=3,\ldots,12$ none lies in either exterior window.

\section{Alternative routes and scope}
The Borcea--Br\"anden calculations certify real-rootedness of 1,242
specific polynomials and 1,097 specific interlacing pairs. However, the
tested global transition operator has 323 negative $2\times2$ minors
among 12,938; an exact witness is $-t$. Thus this total-nonnegativity
route fails, without ruling out every possible stability-preserving
operator argument.

The Lee--Yang substitution
$H(t)=t^mQ(t+t^{-1})$, followed by exact Sturm isolation, certifies
individual supplied reciprocal polynomials. It supplies no uniform
in-$k$ Sturm or discriminant argument. Exact coefficient checks confirm
palindromicity for the computed cases, but do not by themselves imply
that all $3k$ roots are real and negative.

\section{Conclusion and reproducibility}
The cubic recurrence is now structurally proved for every $k\geq3$ by
the exact frontier transfer, with the auxiliary seed specified above.
The signed branch identity at the two $q_6$ contacts, the three-window
phase samples, and the global interval error bound prove exactly $3k$
distinct negative roots for every actual family member $k\geq3$.
Exact Sturm verifies the first actual case $k=2$, completing the
real-rootedness theorem for every $k\geq2$. The auxiliary seed
$\widehat H_{1,4}$ separately has three negative roots. For the
additional finite checks $k=3,\ldots,12$, all roots lie in
$[-79,-1/79]\subset I_4$.
The machine-readable identities and interval data are in
\texttt{certificado\_teorema\_n4.json}. Run the reproducible checks in
this order:
\begin{itemize}
\item \texttt{certify\_n4\_unconditional\_theorem.py}
\item \texttt{verify\_n4\_frontier\_transfer.py}
\item \texttt{audit\_n4\_right\_window.py}
\item \texttt{audit\_n4\_q6\_branch.py}
\end{itemize}
\end{document}
